\documentclass[12pt]{article}
\parindent 0.0in
\parskip 0.1in
\pagestyle{empty}
\def\R{\mbox{I}\negthinspace \mbox{R}}
\topmargin -.5in
\textheight 8.0in
\textwidth 6.0in
\oddsidemargin 0.2in
\usepackage{amsmath}
\usepackage{hyperref}

\begin{document}

\newcommand{\bm}[1]{\mbox{\boldmath$#1$}}

\begin{center}{\bf Homework Set \#1 -- Due 11:59pm, Friday, April 10.} \\ {\sc MATH 447 Spring 2026}\end{center}
\textbf{Note}: 
\begin{itemize}
\item This homework assignment is loosely based on the Chapter 2 and MATH 342 (linear regression course) material. 
\item This homework assignment is designed to make sure you know how to manipulate data using R.
\item For your convenience, all the data files and the .tex file for this problem are included.
\item Your need to submit two files: A PDF file which contains all the answers and an R script file (with an ``.R'' extension) containing all the code.
\item Any answers presented in the R script file only will not be graded.
\item Your answers must be presented in the same order as the problems given in this assignment, including all graphs and {\tt R} outputs, if applicable (i.e., the graphs and outputs must be in the same order and must not be relegated to the back of your assignment). 
\item Use your judgment to include an appropriate amount of {\tt R} code and output in your answers as necessary. For example, it is not necessary for me to see every single observation of the dataset you have used.
\item All the {\tt R} code must be well documented and submitted as a single R script file (which has the ``.R'' extension) to Canvas. 
\end{itemize}

The grading scheme is as follows:
\begin{itemize}
\item This homework assignment will be graded out of $30$ points.
\item Each problem will be graded out of $10$ points.
\item If the R script file is missing or incomplete, at most $5$ points will be deducted.
\end{itemize}

\textbf{Note}: This course assumes that you have at least some knowledge and experience of R programming or something similar. If you are not sure how it works, please see the instructor during the office hours. 

\textbf{Textbook Problems}: Refer to \textit{An Introduction to Statistical Learning with Applications in R}, First Edition by James, Witten, Hastie, and Tibshirani.\\

\begin{enumerate}
\item  \#2.4.2. Note that $n = \mbox{Number of observations}$ and $p = \mbox{Number of explanatory variables}$.\\

\item See S26HW1code.R and answer all the questions from (a) to (i). These questions can be found in the R code, but be sure to present your answers in the PDF file. The dataset is obtained from\\
\scriptsize
\url{https://www.kaggle.com/datasets/farazrahman/earthquake}.\\

\normalsize
\item  \#2.4.8. For (a), do not forget to set {\tt header=TRUE}. Before doing (c) (ii) and (iii), set {\tt college\$Private<-as.factor(college\$Private)}. \\
\end{enumerate}
\end{document}